\section{What partial observations can determine}
\label{l:introduction}

A permutation can look uniform on many labels while retaining information about the whole deck. Parity gives the simplest example: the uniform law on the alternating group has uniform images on any set of at most $n-2$ labels. Thus partial observations require an additional argument before they can establish convergence on $S_n$. This paper develops that argument. Its input is quantitative information about partial permutations; its output is a bound on every nontrivial Fourier matrix, with the sign representation treated separately.

Let $X$ be a set of $n$ labels and $G=\Sym(X)$. We regard $g\in G$ as the map from initial labels to final positions. For $M\subseteq X$, let $H_M=\Sym(M)$ fix $X\setminus M$ pointwise. Two permutations have the same images on $X\setminus M$ precisely when they lie in the same left coset $gH_M$. Consequently an upper bound on the distribution of those images is exactly an upper bound on the masses of these cosets. The elements of $gH_M$ are obtained by multiplication on the right; this convention matters when image observations are compared with inverse-image observations.

For a nonnegative measure $f$ on $G$ and an irreducible unitary representation $\rho:G\to U(V_\rho)$, write
\[
 \widehat f(\rho)=\sum_{g\in G}f(g)\rho(g),\qquad D_\rho=\dim V_\rho.
\]
Our Hilbert--Schmidt norm is unnormalized: $\|A\|_{\HS}^2=\tr(A^*A)$. Let $D_\lambda$ denote the degree of the irreducible of $S_m$ indexed by a partition $\lambda\vdash m$, and set
\[
 C_u=\sup_{m\ge1}\sum_{\lambda\vdash m}D_\lambda^{-u}\quad(u>0).
\]
These constants are finite, and the sums with the row and column partitions removed tend to zero. A short proof is given in Lemma~\ref{l:degree-basic}; Section~\ref{l:degrees} retains the distinct counting proofs that give more detailed information.

\begin{theorem}[Complementary cosets control Fourier matrices]
\label{l:main}
Partition $X$ into $b$ nonempty sets $M_1,\ldots,M_b$. Put $H_i=H_{M_i}$. Suppose $f$ is a nonnegative measure on $G$ of mass at most one and, for every $g\in G$ and $i$,
\begin{equation}\label{l:main-cap}
 f(gH_i)\le \frac{B}{[G:H_i]}.
\end{equation}
For every $u>0$ and every irreducible $\rho$ of degree $D$,
\begin{equation}\label{l:main-hs}
 \|\widehat f(\rho)\|_{\HS}^2
 \le bBC_uD^{-1+(u+2)/b}.
\end{equation}
There is also an independent orbit-span proof of the bound
\begin{equation}\label{l:main-orbit}
 \|\widehat f(\rho)\|_{\op}^2
 \le bBC_uD^{-1+(u+2)/b}.
\end{equation}
No restriction is imposed on block sizes or on the multiplicities in $\rho|_{H_1\times\cdots\times H_b}$.
\end{theorem}

The Hilbert--Schmidt conclusion implies the operator-norm conclusion, but the two proofs use different information. The orbit proof follows one vector and controls the dimension of its orbit under a small subgroup type. The isotypic proof instead bounds the entire Fourier matrix by cosetwise Plancherel. We give the orbit proof first because it isolates the multiplicity issue in a concrete linear map, and then prove the stronger norm conclusion independently.

The principal consequence is a transfer from observations of $7n/8$ labels to the product of just two independent permutations. Suppose $8$ divides $n$, the average total-variation error of these observations tends to zero, and the sign mean is $o(1)$. One can choose eight complementary blocks whose errors sum to $o(1)$, then retain a common subprobability of mass $1-o(1)$ with coset cap one. With $b=8$ and $u=1$, Theorem~\ref{l:main} gives squared Hilbert--Schmidt norm at most $8C_1D^{-5/8}$. Two convolutions have nonsign Plancherel contribution at most $(8C_1)^2\sum D^{-1/4}$, which tends to zero. The discarded mass and sign coefficient also vanish. Corollary~\ref{l:isotypic-completion} gives the precise conclusion. The orbit proof alone gives four convolutions, as recorded in Corollary~\ref{l:orbit-completion}. The distinction is the dimension factor saved by controlling the whole matrix, not a change in the observed marginals.

\subsection{The first proof and its extensions}

Restriction to $H_1\times\cdots\times H_b$ decomposes an irreducible space into tensor-product types, each tensored with its multiplicity space. The product of the factor degrees is at most $D$. Hence some factor has degree at most $D^{1/b}$. Assign the whole type to one such factor. This produces an orthogonal decomposition of any test vector into $b$ pieces. The orbit of one piece under its assigned subgroup has dimension at most the sum of the squares of the small factor degrees. Multiplicity causes no extra factor: on all copies of one type the group acts by the same matrix tensored with identity. The bound $C_u$ now controls that sum. Averaging the projections onto translates of the orbit span, Schur's lemma turns its dimension into a factor $1/D$. This proves \eqref{l:main-orbit}.

The next proof retains more of the Fourier matrix. The central idempotent of a subgroup type projects onto all its copies. Cosetwise convolution bounds the squared Hilbert--Schmidt norm of the projected transform, summed over every ambient irreducible. The small subgroup types cover the original representation; taking a trace against a positive operator completes \eqref{l:main-hs}. This also explains why one must count types once rather than count their multiplicities.

The complete fixed-coset route is proved in Sections~\ref{l:foundations}--\ref{l:core-section}. Sections~\ref{l:comparisons-section} and~\ref{l:random-partition-section} compare alternative coefficient arguments and retained measures for the same marginal information. Section~\ref{l:degrees} develops the shape-sensitive degree estimates needed later, and Section~\ref{l:characters-section} constructs one-dimensional characters in subgroups of controlled index. Supplementary tableau and hook proofs are grouped in Appendix~\ref{l:degree-supplement}. Appendix~\ref{l:transfer-supplement} retains the projection-averaging and longest-line peeling methods, together with explicit parameter substitutions. These independent routes remain available without interrupting the main transfers.

The remaining results distinguish three kinds of additional information. First, image and inverse-image caps control opposite sides of a Fourier matrix. Section~\ref{l:two-sided-section} places those sides together in a product of two laws. Section~\ref{l:equivariant-section} instead works with transition kernels: the sign contribution is then an operator on a multiplicity space, not merely the scalar sign mean of a group law.

Second, the quantifiers on an information bound determine which changes of law it permits. Section~\ref{l:domains-section} constructs retained laws by selecting good omitted sets separately for each permutation, under the same averaged marginal hypothesis as the fixed-partition route. Section~\ref{l:tilts-section} assumes an entropy comparison for every law on one common event; it can therefore condition on a rare representation-coefficient event. Section~\ref{l:palettes} obtains analogous rare-event control from a pointwise bound on accumulated reveal errors. In contrast, the single-law entropy bounds of Section~\ref{l:replica-section} give a growing coset cap by clipping. They control only sufficiently large dimensions; Section~\ref{l:forests} supplies a precise hybrid conclusion when a separate law handles the remaining diagrams by long-row decay and exact vanishing on diagrams with more than $n/2$ rows.

Finally, Section~\ref{l:sec:tabloid} uses a conditional product structure: after fixing an environment, the two half-permutations are independent. A signed tabloid basis turns their image caps into matrix bounds, while a preceding law's inverse-image caps average positive operators. Appendix~\ref{l:tabloid-supplement} gives an independent arm--leg construction of the signed representation used there. These different hypotheses are kept explicit; no mixing conclusion is inferred by substituting one kind of partial information for another.

\subsection{Relation to earlier methods}

Fourier analysis of random walks on finite groups converts distributional convergence into estimates of matrix transforms. Diaconis and Shahshahani's analysis of random transpositions is a fundamental example \cite{diaconis-shahshahani}; the Plancherel and total-variation steps used here belong to that framework. We keep matrix norms throughout because the laws considered here need not be central and their transforms need not be normal.

The representation-theoretic ingredients are classical. Specht modules, Young branching, induced representations from Young subgroups, and the Littlewood--Richardson rule supply the restriction and signed-tabloid constructions; see James \cite{James1978} and Sagan \cite{sagan}. The hook formula of Frame, Robinson, and Thrall \cite{frame-robinson-thrall} and elementary standard-tableau constructions supply degree lower bounds. The reciprocal-degree limits are part of the symmetric-group Witten-zeta theory; see Liebeck and Shalev \cite[Proposition~2.5 and Theorem~2.6]{liebeck-shalev2004}. We give elementary proofs suited to the different transfer arguments and retain their more detailed tableau bounds. The orbit-span and coset estimates are proved explicitly. Relative entropy enters through its chain rule and its variational behavior under conditioning; the rare-event arguments specify the changed law explicitly.

One motivation is the Thorp shuffle, introduced in \cite{Thorp1973}.
Entropy methods yielded an $O((\log n)^4)$ mixing bound for even decks
and an $O((\log n)^3)$ bound for power-of-two decks
\cite{Morris2009,Morris2013}. These bounds count elementary
cut-and-interleave shuffles; on $n=2^d$ cards, a complete coordinate
sweep consists of $d$ such steps. The entropy arguments reveal and
compare card trajectories, providing a methodological antecedent to
the partial-information hypotheses considered here.

More directly, Czumaj and V\"ocking studied the joint randomization
of a fixed fraction of labeled cards under the Thorp shuffle
\cite{CzumajVocking2014}. Czumaj subsequently obtained full-permutation
sampling by recursively applying partial-permutation randomizers on
successively smaller sets \cite{Czumaj2015}. The present transfer has
a different form: complementary quotient estimates for a single law
control its noncentral Fourier matrices and, with separate sign
control, imply mixing after a fixed number of independent compositions
of that same law. No shuffle construction or physical-time bound is
assumed in the abstract proofs.

Three companion papers construct inputs for the later transfers.
The coordinate-frame paper proves the all-tilt entropy estimate in
Lemma~12.13 and establishes the quotient and retained sign-multiplicity
estimates within the proof of Theorem~14.2 \cite{companion-frames}.
The conditional-information paper supplies the fixed-list estimate in
Proposition~3.3; the palette estimate and common-event construction in
Proposition~13.2 and equations~(13.2)--(13.3); the random-domain
estimate in Theorem~12.2; the weak-entropy and sparse-Fourier inputs in
Propositions~15.1 and~15.4; and the entropy assertion~(16.1) of
Theorem~16.1 \cite{companion-information}.
The conditional-permutation paper constructs the two-half environment
and its independent internal permutations in the proof of Theorem~7.1
and Lemma~7.2 \cite{companion-kernels}.
These estimates are established before the cited papers apply the
transfer results here. Their full-deck mixing conclusions are
applications, not hypotheses of the present proofs; their discarded-mass,
parity, and physical-time checks remain part of those applications.
